\subsection{Query Language}
The basic structure of a \sql\ statement is \texttt{SELECT I FROM T WHERE C},
which corresponds to $\Pi_{\texttt{I}}(\sigma_{\texttt{C}} \texttt{T})$.

We can set \texttt{I = *} if we want to return all columns for a table, or we can even set \texttt{I = A.*}, for a table \texttt{A} in \texttt{T} called \texttt{A},
i.e. with e.g. \texttt{T = Table A}.
To rename, we can set \texttt{I = Column AS Name, Column2 AS Name2}, etc

\inlinetheorem Every SPJR \acrshort{ra} expression can be written in \texttt{SELECT ... FROM ... WHERE ...} form:
\begin{tables}{lll}{Operation   & Notation                   & SQL}
              Selection         & $\sigma_c(R)$              & \texttt{SELECT * FROM R WHERE c;}              \\
              Projection        & $\Pi_{A_1, \ldots, A_n} R$ & \texttt{SELECT A1, \ldots, An FROM R;}         \\
              Cartesian Product & $R_1 \times R_2$           & \texttt{SELECT * FROM (R1, R2);}               \\
              Rename            & $\rho_{a, b, c} R$         & \texttt{SELECT A as a, \ldots, B as c FROM R;} \\
              Union             & $R_1 \cup R_2$             & \texttt{R1 UNION R2;}                          \\
              Difference        & $R_1 - R_2$                & \texttt{R1 EXCEPT R2;}                         \\
              Intersection      & $R_1 \cap R_2$             & \texttt{R1 INTERSECT R2;}                      \\
\end{tables}

Since \sql\ implements \gls{bag} and not set semantics, there is a \texttt{DISTINCT} keyword, which is used to remove duplicates.
It is to be applied in the \texttt{SELECT} clause (\texttt{I} above), e.g.
\mint{sql}|SELECT DISTINCT name FROM data_table;|

For the last three operations in the table above, \texttt{R1} and \texttt{R2} typically are other \sql\ statements, example:
\mint{sql}|(SELECT name FROM FirstTable) UNION (SELECT name FROM SecondTable)|
In addition, to apply \textit{bag semantics}, as opposed to \textit{set semantics} append \texttt{ALL} to the \texttt{UNION} expression.

We can also name the tables, e.g. $\texttt{T} = \texttt{One o, Two t}$, and then access columns from a specific table using
$\texttt{I} = \texttt{o.Col, b.Col}$, or the like.

String concatenation works using \texttt{CONCAT('string', 'string', 'string', ...)}, or using \texttt{'string' || 'string' || \dots}.
Note that SQL uses single quotes for Strings, and double quotes for renaming columns (using \texttt{AS} statements, or in \texttt{SELECT})

Dates work as you'd expect. We can create a new date using \texttt{Date('ISO-date-string')}. \texttt{DATETIME} combines date and time and \texttt{TIME} is just the time.
To compute time delta, we can use \texttt{DATEDIFF('interval', DateOne, DateTwo)}, where the interval can be:
\texttt{year, quarter, month, dayofyear, day, week, weekday, hour, minute, second, millisecond}.
To get just the year from an existing date object (or string), use \texttt{YEAR}. The result is an integer and thus you can use comparison operations with numbers.

The current date, time, year, etc is provided using \texttt{current\_date} (etc).

There are also conditional statements, namely \texttt{CASE WHEN condition1 THEN result1 WHEN condition2 THEN result2 ELSE else END}, or of course, more conditions.
This can even be used inside aggregations, etc!
